%!TEX program = lualatex

\documentclass[portuguese, 11pt, a4paper, titlepage, oneside]{article}

\usepackage{geometry}
\usepackage[T1]{fontenc}
\usepackage{graphicx}
\usepackage{parskip}
\usepackage{amssymb}
\usepackage[portuguese]{babel}
\usepackage[framed, numbered]{matlab-prettifier}
\usepackage{subcaption}
\usepackage{wrapfig}
\usepackage{hyperref}

% Title page config
\title{
  \includegraphics[width=10cm]{assets/deti_horizontal_preto.pdf}
  \vskip 2em
  \huge Aspetos Profissionais e Sociais\\
  da Engenharia Informática
  \vskip 1em
	{\bfseries Trabalho individual 1}
}

\author{Rúben Garrido\\NMec: 107927}
\date{\today}
% Begin document
\begin{document}

% Title
\maketitle

\section{Utilização do ChatGPT pelos alunos, no contexto da UC de APSEI}
A inteligência artificial baseada em \textit{Large Language Models} (neste caso, o ChatGPT) revolucionou a forma como os humanos pensam e interagem com a tecnologia. Todas as áreas foram influenciadas, embora a educação tenha tido um dos maiores destaques.

Esta mudança de paradigma deve-se sobretudo à utilização do ChatGPT (bem como de outros LLMs) pelos alunos, por forma a "facilitar" a curva de pensamento necessária para realizar um determinado trabalho. No entanto, o uso excessivo e irresponsável da inteligência artificial levou ao surgimento de respostas meramente criadas por esta, com imprecisões e padrões repetidos.

\section{O que falta abordar}
\begin{itemize}
  \item If yes, should students be expected to use ChatGPT?
  \item If yes, should students be incentivized to use ChatGPT in which way?
  \item If no, why not?
  \item And can/should professors follow the same rules as the students?
\end{itemize}
\end{document}